
\def\PUDcodigo{ECA205}

\def\PUDcodacad{04507.34}

\def\PUDdisciplina{Processamento Digital de Sinais}

\def\PUDsemestre{S7}

\def\PUDhoras{80}

%NOVOS ---------------------------------------------------
\def\PUDhorasTeoria{64}
\def\PUDhorasPratica{16}

\def\PUDinterECA{ \hyperref[PUD:ECA211]{Eletrônica I (04507.15)}

\hyperref[PUD:ECA220]{Sistemas Lineares (04507.26)} 

\hyperref[PUD:ECA202]{Controle II (04507.40)}
}

\def\PUDatuaisECA{---}

\def\PUDrecursos{Projetor multimídia; Quadro branco e pincel; Aulas em laboratório de informática.}

%----------------------------------------------------------

\def\PUDcreditos{4}

\def\PUDprerequisito{ \hyperref[PUD:ECA209]{ECA209} }

\def\PUDpreacad{ \hyperref[PUD:ECA220]{Sistemas Lineares (04507.26)} }

\def\PUDobjetivos{Compreender aspectos teóricos e práticos relativos ao processamento de sinais discretos.  Aplicar técnicas para aquisição, filtragem e análise de sinais por meio de computador.}

\def\PUDementa{Introdução aos sinais e sistemas discretos.  Sinais e sistemas discretos.  Transformada Z.  Amostragem de sinais contínuos no tempo.  Análise de sistemas lineares e invariantes.  Estruturas de sistemas discretos. 
Técnicas e projetos de filtros.  Transformada discreta de Fourier.  Algoritmos rápidos para a transformada de Fourier.  Projeto de filtros digitais.  Simulações de filtros digitais.}

\def\PUDconteudo{ & Introdução aos sinais e sistemas discretos: representação matemática de sinais contínuos e discretos, sinais periódicos e aperiódicos, sinais contínuos e discretos básicos, operações sobre sinais discretos, convolução, propriedades de sistemas.  Exemplos. 
\\ 
 & Transformada Z: definição da transformada Z, pólos e zeros, região de convergência e transformada inversa, propriedades da transformada, solução de equações a diferenças com coeficientes constantes.  Exemplos. 
\\ 
 & Amostragem de sinais contínuos no tempo: representação de um sinal contínuo no tempo pelas suas amostras, amostragem por trem de impulsos, teorema da amostragem, reconstrução de um sinal contínuo no tempo a partir de suas amostras, sub-mostragem e aliasing.  Exemplos. 
\\ 
 & Análise de sistemas lineares e invariantes: resposta em freqüência de sistemas LTIs;  sistemas caracterizados por equações de diferença com coeficientes constantes;  resposta em freqüência de sistemas caracterizados por funções racionais;  relações entre magnitude e fase;  sistemas passa-tudo, de mínima fase e de fase linear.  Exemplos. 
\\ 
 & Estruturas de sistemas discretos: representação em diagrama de blocos de equações de diferença com coeficientes constantes;  estruturas básicas de sistemas IIR;  formas transpostas;  estruturas básicas de redes para sistemas FIR;  efeitos da precisão numérica finita e da quantização;  propagação do ruído em filtros digitais;  análise de ponto-fixo e ponto-flutuante em projetos de filtros digitais.  Exemplos. 
\\ 
 & Projeto de filtros digitais de sinais: filtros IIR e FIR, projeto de filtros digitais IIR a partir de filtros analógicos, transformação bilinear, propriedades dos filtros FIR, projetos de filtros FIR usando janelas, comparação de filtros analógicos e filtros digitais, projeto de filtros com aplicação na redução de ruído em sinais.  Exemplos. 
\\ 
 & Transformada de Fourier discreta: sinais periódicos e sua representação pela série discreta, representação de sequências de duração finita pela transformada de Fourier, convergência, propriedades da transformada de Fourier no tempo discreto, transformada inversa, sistemas lineares descritos por equações a diferenças de coeficientes constantes, aplicações.  Exemplos. }

\def\PUDmetodo{Aulas expositórias que podem ser teóricas e/ou práticas, onde as práticas no laboratório serão marcadas no decorrer da disciplina.}

\def\PUDavalia{A avaliação será feita com aplicação de provas teóricas e/ou práticas, além da possibilidade de inclusão de trabalhos e seminários no decorrer da disciplina.}

\def\PUDbibbasica{ & \href{www.biblioteca.ifce.edu.br/index.asp?codigo_sophia=40232}{Oppenheim}, Alan V.; Schafer, Ronald W.  Processamento em tempo discreto de sinais.  3º ed.  São Paulo, SP: Pearson Education do Brasil, 2012.  665p.  ISBN: 9788581431024.
\\ 
 &  \href{www.biblioteca.ifce.edu.br/index.asp?codigo_sophia=39837}{Lathi}, B. P.  Sinais e sistemas lineares.  2º ed.  Porto Alegre, RS: Bookman, 2007.  856p.  ISBN: 9788560031139.
\\ 
&  \href{www.biblioteca.ifce.edu.br/index.asp?codigo_sophia=63268}{Roberts}, Michael J.  Fundamentos em sinais e sistemas.  São Paulo, SP: McGraw-Hill, 2009.  764p.  ISBN: 9788577260386. }
%&  HAYKIN, Simon;  VEEN, Barry Van.  Sinais e sistemas.  Porto Alegre (RS): Bookman, 2001/2007.  668p.   ISBN 9788573077414 .}

\def\PUDbibcomplementar{ & \href{www.biblioteca.ifce.edu.br/index.asp?codigo_sophia=39861}{Oppenheim}, Alan V.  Sinais e sistemas.  2º ed.  São Paulo, SP: Pearson Prentice Hall, 2010.  568p.  ISBN: 9788576055044.
%&  KUO, Sen M.  (Sen-Maw);  GAN, Woon-Seng.  Digital signal processors: architectures, implementations and applications.  Upper Saddle River (NJ): Pearson Prentice Hall, 2005.  601 p.   ISBN %9780471295464
\\ 
 &  \href{www.biblioteca.ifce.edu.br/index.asp?codigo_sophia=62379}{Girod}, Bernard; Rabenstein, Rudolf; Stenger, Alexander.  Sinais e sistemas.  Rio de Janeiro, RJ: LTC, 2003.  340p.  ISBN: 8521613644.
%&  PROAKIS, John G.;  MANOLAKIS, Dimitris G.  Digital signal processing.  4. ed.  Upper Saddle River (NJ): Pearson Education, 2007.  948 p.  ISBN 9780131873742 
\\ 
 &  \href{www.biblioteca.ifce.edu.br/index.asp?codigo_sophia=24208}{Stallings}, William.  Redes e sistemas de comunicação de dados : teoria e aplicações corporativas.  Rio de Janeiro, RJ : Elsevier, 2005.  449p.  ISBN: 8535217312.
\\
 &  \href{www.biblioteca.ifce.edu.br/index.asp?codigo_sophia=55613}{Ogata}, Katsuhiko.  Engenharia de Controle Moderno.  E-book.  4º ed.  Pearson.  800p.  ISBN: 9788587918239.
\\
 &  \href{www.biblioteca.ifce.edu.br/index.asp?codigo_sophia=41966}{Boyce}, William E.  Equações diferenciais elementares e problemas de valores de contorno.  9º ed.  Rio de Janeiro, RJ: LTC, 2013.  607p.  ISBN: 9788521617563. }
%  PROAKIS, John G.;  MANOLAKIS, Dimitris G.  Digital signal processing.  4. ed.  Upper Saddle River (NJ): Pearson Education, 2007.  948 p.  ISBN 9780131873742.

\def\PUDautor{Geraldo Luis Bezerra Ramalho}

\def\PUDdata{2013-05-20}

\def\PUDrevisao{3}

\def\PUDrevisor{Adriano Holanda Pereira}

\def\PUDdatarevisao{2019-05-15}

\PUDcorpo
